% Chapter 2 can be included from a larger thesis with \input or \include.
% Existing citation keys follow the user-maintained references.bib file.
% References first introduced for this chapter are grouped in a clearly marked
% block at the end of references.bib.

\chapter{Literature Review}
\label{chap:literature_review}

\section{Introduction}
\label{sec:lit_introduction}

Music-driven humanoid motion lies at the intersection of music information
retrieval, data-driven character animation, real-time control, and humanoid
robotics. These fields share an interest in temporal behaviour, but they make
different assumptions about what information is available and what constitutes
a successful output. Music information retrieval is primarily concerned with
extracting useful structure from sound. Dance-generation research usually aims
to synthesise plausible human motion from a complete music clip. Interactive
animation instead prioritises responsiveness and continuity under online user
control. Humanoid robotics adds embodiment, contact, and feasibility constraints
that are absent from an unconstrained virtual character.

This chapter reviews these areas from the perspective of the problem defined in
the preceding chapter: causal selection and continuous playback of robot-ready
dance motions from streaming audio. The review distinguishes three
forms of music-driven motion system. An \emph{offline generative system} may
observe a complete audio sequence and synthesise a complete motion sequence. A
\emph{streaming generative system} produces new motion while restricting its
access to future audio. A \emph{retrieval-and-control system} selects and adapts
motion from a finite database. The proposed Humanoid Matcher belongs to the
third category. Its closest conceptual predecessors therefore include not only
music-to-dance models, but also motion graphs, motion matching, beat tracking,
and humanoid motion retargeting.

The purpose of the review is not to claim that retrieval is a substitute for
generative modelling. Instead, it identifies the trade-offs among novelty,
latency, interpretability, motion provenance, and embodiment safety. These
trade-offs motivate a hybrid architecture in which audio analysis, cross-modal
retrieval, temporal decision rules, phase control, and prevalidated robot
motion each address a different part of the overall problem.

\section{Music Representation for Motion Control}
\label{sec:lit_music_representation}

\subsection{Temporal and semantic information in music}
\label{subsec:lit_temporal_semantic}

A music representation for dance control must describe both \emph{when} motion
events should occur and \emph{what kind} of movement may be appropriate.
Rhythmic descriptors such as tempo, onset strength, beat confidence, and
inter-beat interval provide temporal cues. Spectral shape, timbre, harmony,
loudness, and learned style embeddings provide broader evidence about musical
content and activity. These two groups are related but not interchangeable. Two
tracks can have similar tempi yet differ substantially in energy and timbre;
conversely, two stylistically similar tracks can require different playback
rates.

Classical music information retrieval systems construct descriptors from
short-time signal analysis. A frame-level descriptor sequence may be written as

\begin{equation}
    \mathbf{f}_{1:T}
    = \mathcal{F}\!\left(\mathbf{a}_{1:S}\right),
\end{equation}

where $\mathbf{a}_{1:S}$ is an audio waveform and $\mathcal{F}$ may include a
short-time Fourier transform, mel-frequency bands, spectral statistics,
chroma, onset strength, or rhythm summaries. Libraries such as Essentia collect
spectral, temporal, tonal, and high-level descriptors in a common analysis
framework \cite{bogdanov2013essentia}. Such descriptors are computationally
attractive because their meaning and cost are explicit. They can also be
normalised independently, which is valuable when microphone gain, playback
volume, or recording conditions change.

However, hand-designed features expose an important design tension. Global
summary statistics are cheap and stable, but they can discard local structure.
Frame-level features retain that structure, but direct comparison becomes more
expensive and sensitive to alignment. The appropriate representation therefore
depends on the decision timescale. A short rolling window is useful for
responsiveness, whereas a longer context is useful for estimating stable style
or tempo. A real-time system commonly needs both: local events for phase
control, and a slower descriptor for motion selection.

\subsection{Learned music embeddings}
\label{subsec:lit_learned_embeddings}

Learned embeddings replace or complement manually chosen descriptors by mapping
audio into a feature space whose geometry is shaped by a training objective.
Music representations can be supervised by genre and mood labels, by
co-occurrence information, or by metadata-derived similarity. In particular,
Alonso-Jim\'enez et al. use associations extracted from Discogs editorial
metadata as contrastive supervision and show that large, weakly labelled music
collections can produce useful general-purpose representations
\cite{alonso2022music}. This line of work is relevant to the Discogs EffNet
backend used by the current Humanoid Matcher catalogue: a learned embedding
can encode stylistic relations that are difficult to state as a small set of
DSP rules.

A learned representation nevertheless introduces new sources of uncertainty.
Its similarity metric reflects the pretraining data and objective rather than
the requirements of robot motion. A track-level style embedding may rank two
pieces of music as similar even when their beat rates would require an
infeasible change in motion speed. Domain shift is also possible when a model
trained on mastered recordings receives reverberant or clipped microphone
audio. For this reason, semantic embedding similarity should not be treated as
a complete cross-modal compatibility function. A hybrid score can retain the
coverage of learned features while applying explicit tempo, rhythm, and
activity terms whose contribution can be inspected and ablated.

Let $\mathbf{e}_q$ and $\mathbf{e}_i$ be standardised query and catalogue
embeddings. A typical content similarity is their cosine similarity,

\begin{equation}
    s_{\mathrm{emb}}(q,i)
    = \frac{\mathbf{e}_q^{\mathsf{T}}\mathbf{e}_i}
           {\lVert\mathbf{e}_q\rVert_2\lVert\mathbf{e}_i\rVert_2}.
\end{equation}

For motion control, this term is more naturally interpreted as evidence for a
related musical neighbourhood than as direct evidence that a particular
trajectory is executable. The distinction motivates hierarchical retrieval:
first identify relevant music, then rank the motions associated with that music
using motion-specific compatibility.

\subsection{Beat and onset analysis under causal constraints}
\label{subsec:lit_causal_beats}

Beat tracking illustrates the difference between analysing music and
controlling motion from music. Ellis formulates beat tracking as a compromise
between local onset evidence and a transition preference for a regular
inter-beat interval \cite{ellis2007beat}. This formulation captures a general
principle: a beat is not merely a spectral peak, but a temporally coherent
hypothesis supported by recent events. For offline analysis, dynamic programming
can optimise a sequence of beat times over a complete region. A live controller
cannot freely revise past beat decisions after the corresponding robot pose has
already been emitted.

Causal beat handling therefore requires explicit confidence and latency
policies. Strong percussive subdivisions can produce twice the intended control
rate, while weak passages can produce missing or unstable events. Accepting
every detected event risks aligning motion key poses to hi-hat subdivisions
rather than perceptually strong beats. Rejecting too many events causes stale
tempo estimates and slow phase correction. Relative onset or loudness contrast
over recent beats offers a causal way to prefer stronger events, but its quality
depends on the length and diversity of the available history.

The important consequence for the present research is that audio retrieval and
beat-phase control should use related but separate signals. Retrieval benefits
from several seconds of context and should not change at every onset. Phase
control requires low-latency event timing but should not reinterpret every
event as a change of musical identity. This separation reduces the risk that a
single noisy audio block simultaneously changes the selected dance and its
playback phase.

\section{Dance Data and Motion Representation}
\label{sec:lit_data_motion}

\subsection{Paired music--dance datasets}
\label{subsec:lit_datasets}

Progress in music-conditioned dance generation has been strongly influenced by
the scale and representation of available paired datasets. AIST++ contains
1,408 three-dimensional dance sequences across ten genres and is accompanied by
aligned music and camera information \cite{li2021ai}. Its size, genre labels,
and synchronisation make it a common benchmark for music-conditioned generation
and a practical source for the catalogue in this dissertation. AIST++ is
especially relevant to retrieval-based work because multiple dance sequences
can be associated with related musical material, permitting the catalogue to
represent one-to-many music--motion relationships.

FineDance extends the emphasis from broad coverage to fine-grained, full-body
choreography. It increases stylistic diversity and includes more detailed hand
movement and genre coverage than earlier benchmarks \cite{li2023finedance}.
The accompanying method also uses retrieval as contextual assistance for
generation, illustrating that retrieval and generative modelling need not be
mutually exclusive. Retrieval can supply long-range examples or style context,
while a generator produces local motion.

Dataset scale does not remove ambiguity from evaluation. Music and dance have
a many-to-many relationship, so a generated or retrieved movement can be
appropriate without matching a held-out reference pose. Dataset labels may
also reflect dance genre, music genre, performer, or recording session in
different proportions. Models can consequently exploit correlations that do
not transfer to live music. Evaluations based only on reconstruction distance
are therefore insufficient; rhythmic alignment, distributional quality,
diversity, transition continuity, and user judgement provide complementary
evidence.

\subsection{Human-body and robot-motion representations}
\label{subsec:lit_motion_representations}

Human dance datasets frequently represent motion through skeletal joint
positions, joint rotations, or a parametric body model. SMPL defines a skinned
human model with pose and shape parameters and is compatible with standard
graphics pipelines \cite{loper2023smpl}. It provides a useful canonical human
representation because rotations preserve articulated structure and the model
can generate consistent joint and surface geometry across performers.

For a target humanoid, however, an SMPL pose is not a command. A robot has a
different kinematic tree, joint axes, degrees of freedom, proportions, and
collision geometry. A representation suitable for visual generation may also
omit velocities, contact states, actuator limits, or world-frame root
requirements. Thus, the apparent output type ``3D motion'' conceals a major
difference between human animation and robot execution.

The representation selected for a catalogue also affects matching and
transition costs. Raw joint angles allow direct comparison in the target
configuration space but require meaningful normalisation across joints.
Cartesian joint positions make geometric similarity clearer but can hide joint
limit differences and amplify root-alignment errors. Velocities reveal whether
two poses are moving toward compatible future states. Foot contacts encode a
discrete support condition that neither pose nor velocity alone captures.
Consequently, a robot transition should normally consider a structured state
rather than a single pose vector.

\section{Music-Conditioned Dance Generation}
\label{sec:lit_dance_generation}

\subsection{Autoregressive and discrete motion models}
\label{subsec:lit_autoregressive}

FACT couples a full-attention cross-modal transformer with the AIST++ dataset
to model long music and motion sequences \cite{li2021ai}. Its transformer-based
formulation demonstrates that temporal context and cross-modal attention can
capture dependencies beyond a local beat. At the same time, autoregressive
generation exposes familiar difficulties: errors can accumulate over long
horizons, and a model conditioned on previously generated frames can become
sensitive to its own imperfect history.

Bailando addresses choreographic structure through a learned discrete motion
codebook and an actor--critic GPT framework \cite{siyao2022bailando}. The
codebook turns continuous pose sequences into motion tokens, while the reward
design encourages beat alignment. Discrete representations can make long-range
modelling more manageable and can preserve recurring movement patterns.
However, token boundaries do not automatically correspond to dynamically safe
robot transition points, and decoding a human pose sequence still leaves an
embodiment conversion problem.

These methods establish two important ideas for the Humanoid Matcher. First,
musical context beyond the instantaneous beat is necessary for coherent
selection. Second, repeated motion patterns are not inherently undesirable;
dance contains motifs whose recurrence supports style and structure. A
retrieval system can exploit this property directly, although its available
motifs remain limited to the catalogue.

\subsection{Diffusion-based dance synthesis}
\label{subsec:lit_diffusion}

Diffusion models treat motion generation as iterative denoising conditioned on
audio and other signals. Listen, Denoise, Action! applies this formulation to
audio-driven motion and supports probabilistic sampling rather than a single
deterministic output \cite{alexanderson2023listen}. EDGE uses a transformer
diffusion model for music-conditioned dance and demonstrates both high-quality
generation and editing operations such as motion in-betweening
\cite{tseng2023edge}. Beat-It introduces explicit beat and key-pose conditions,
showing the value of representing rhythmic constraints directly rather than
expecting them to emerge only from a global audio embedding
\cite{huang2024beat}.

The principal attraction of diffusion is its capacity to represent multiple
valid motions for the same music. This is well matched to the multimodal nature
of dance. Iterative sampling can also integrate diverse constraints. The cost
is that denoising introduces computation and, in a conventional formulation,
often assumes a fixed sequence whose future frames are revised together. A
motion that looks coherent after full-sequence denoising does not by itself
establish causal, bounded-latency behaviour.

For this reason, offline diffusion results and real-time control results should
not be compared solely through motion quality metrics. They solve different
information problems. If an offline model observes future musical events, it
can anticipate an upcoming accent; a causal controller can respond only after
sufficient evidence has arrived. The fair comparison must include latency,
look-ahead, revision policy, and how committed motion history is handled.

\subsection{Retrieval within generative systems}
\label{subsec:lit_retrieval_generation}

FineDance incorporates a retrieval mechanism to improve genre consistency and
long-term coherence before or alongside generation \cite{li2023finedance}.
This hybrid design points to a broader interpretation of retrieval. A database
need not merely return a final clip; it can provide contextual priors, examples,
or candidates that are subsequently adapted. The quality of retrieval then
depends on the correspondence between the query feature space and the desired
motion property.

In the Humanoid Matcher, retrieval has a stricter role: the selected item is a
prevalidated robot trajectory rather than a context example for a generator.
This reduces representational freedom but strengthens motion provenance. Every
candidate can be inspected offline, linked to its source audio and retargeting
metadata, and rejected if it fails a preflight check. The comparison suggests a
continuum rather than a binary distinction. At one end, pure generation offers
high novelty with less explicit provenance. At the other, pure retrieval offers
strong provenance with finite coverage. Future systems may use retrieval to
propose safe motion primitives and generation to adapt them within controlled
bounds.

\section{Real-Time Audio-Driven Character and Humanoid Control}
\label{sec:lit_realtime_control}

\subsection{Causality, latency, and committed history}
\label{subsec:lit_causality}

For an online system, causality can be expressed as

\begin{equation}
    \mathbf{q}_t
    = \mathcal{G}\!\left(\mathcal{A}_{\leq t-\delta},
                         \mathbf{q}_{<t}\right),
\end{equation}

where $\delta$ represents acquisition and processing delay. A strictly causal
system has no access to $\mathcal{A}_{>t}$, and a bounded-latency system must
also place a practical upper bound on $\delta$. In addition, previously emitted
poses $\mathbf{q}_{<t}$ are committed: an algorithm may update its internal
belief about the music, but it cannot alter the robot's physical past.

This formulation reveals three distinct latency sources. \emph{Perceptual
latency} comes from the audio context needed to estimate a reliable feature.
\emph{Computational latency} comes from feature extraction, inference, search,
and planning. \emph{Scheduling latency} comes from deliberately waiting for a
musically or mechanically suitable switch boundary. Reducing only model
inference time does not guarantee a responsive system if a six-second feature
window is recomputed infrequently or if every change waits several bars.

\subsection{Direct audio-to-humanoid policies}
\label{subsec:lit_freestyle}

\emph{Do You Have Freestyle?} introduces RoboPerform, an audio-conditioned
humanoid locomotion framework that aligns audio and motion representations and
distils a mixture-of-experts teacher into a diffusion-based student policy
\cite{li2026you}. A central contribution is retargeting-free control: the
student produces humanoid actions from audio rather than first generating an
explicit human sequence and passing it through a separate retargeter. This can
reduce accumulated error and latency across pipeline stages. The work also
demonstrates expressive behaviour on a Unitree G1 platform, making it more
directly relevant to humanoid performance than animation-only dance models.

RoboPerform and the Humanoid Matcher address related goals through different
assumptions. RoboPerform learns an action policy and can synthesise behaviour
beyond a fixed set of stored trajectories. The matcher assumes that a library
of actions has already been retargeted and checked. The former places more of
the problem in representation learning and policy training; the latter places
more of it in catalogue construction, similarity design, selection stability,
and transition planning. The contrast is useful because it identifies two
different ways of handling embodiment: learn it inside the policy, or validate
it before retrieval.

\subsection{Streaming diffusion forcing}
\label{subsec:lit_discoforcing}

DiscoForcing explicitly formulates audio-driven character motion as strict
causal streaming with bounded latency \cite{ji2026discoforcing}. It combines a
causal music encoder with diffusion forcing, a hybrid temporal training
schedule, and history-guided sampling. The model is designed to remain
responsive to new audio while preventing long-term degradation from stale or
imperfect generated history. Its demonstrations span animated characters and
humanoid deployment, illustrating that streaming generation can be treated as
a first-class problem rather than an offline model run on shorter windows.

The importance of DiscoForcing to this dissertation is methodological. It
provides a stronger standard for the term ``real time'': causal access,
specified latency, continuous history, and online output should all be stated.
The Humanoid Matcher does not use diffusion forcing, but it faces analogous
stability questions. A generative model must decide how much noisy history to
condition on; a retrieval system must decide how much ranking history is enough
to justify a switch. Both systems balance responsiveness against temporal
consistency, although they implement that balance at different levels.

There is also an important difference in failure behaviour. A streaming
generator may gradually drift or produce locally implausible motion, whereas a
retrieval system is bounded by the quality of its stored clips but can still
fail through a poor match, phase error, or discontinuous transition. Therefore,
the evaluation metrics should differ. Generative diversity and distributional
quality are central for DiscoForcing; catalogue recall, switching stability,
entry-state cost, and deadline misses are central for the Humanoid Matcher.

\section{Example-Based Motion Synthesis and Motion Matching}
\label{sec:lit_motion_matching}

\subsection{Motion graphs and database reuse}
\label{subsec:lit_motion_graphs}

Long before modern audio-conditioned diffusion, data-driven animation
investigated how captured motions could be rearranged without losing their
realism. Motion graphs construct a directed graph whose edges contain captured
motion or generated transitions and whose walks form new sequences
\cite{kovar2023motion}. The critical operation is the discovery of states
at which one motion can transition into another with limited visual
discontinuity. This reframes motion synthesis as search over reusable examples.

Motion graphs are conceptually relevant to a finite dance library, but a dance
matcher introduces additional structure. Candidate transitions are not chosen
only to follow a user-specified spatial path; they are conditioned by music,
beat boundaries, motion salience, and humanoid support state. Moreover, a
transition graph built from every compatible frame pair may become large. The
present system instead performs a local entry search when a motion switch is
accepted. This avoids storing a dense graph while retaining the central insight
that database clips should be joined at compatible states rather than arbitrary
timestamps.

\subsection{Motion matching and learned compression}
\label{subsec:lit_learned_motion_matching}

Motion matching searches a motion database at runtime for a frame whose pose
and future trajectory best fit the current character state and desired control.
Learned Motion Matching compresses this behaviour into neural networks while
retaining the responsive, data-driven character of database search
\cite{holden2020learned}. The work highlights the main strengths of example-based
control: captured detail can be preserved, behaviour is guided by an explicit
query, and selection can respond continuously to an external target.

The Humanoid Matcher uses the term \emph{matcher} in a related but more
hierarchical sense. Its primary query is music rather than a desired locomotion
trajectory, and it initially ranks complete motion candidates associated with
catalogue tracks. Only after a candidate is selected does it search entry
states using the current robot pose and velocity. This division addresses a
scaling and semantics problem: joint-state proximity alone cannot establish
that a dance matches the current music, while music similarity alone cannot
establish that the transition is continuous.

A general hierarchical score can be expressed as

\begin{equation}
    i^{*}
    = \arg\max_{i \in \mathcal{C}(q)}
      \left[
        w_m s_{\mathrm{music}}(q,i)
        + w_{\tau}s_{\mathrm{tempo}}(q,i)
        + w_k s_{\mathrm{key}}(q,i)
        + w_a s_{\mathrm{activity}}(q,i)
      \right],
\end{equation}

where $\mathcal{C}(q)$ is a candidate set produced by track-level retrieval.
This formulation makes the compatibility dimensions visible. It also supports
ablations that would be harder to interpret if the entire cross-modal relation
were represented by one end-to-end latent score.

\subsection{Selection stability and diversity}
\label{subsec:lit_selection_stability}

Nearest-neighbour retrieval is stateless: a small score change can replace the
current winner. An online controller is stateful, and switching itself has a
cost. When two candidates have nearly equal scores, immediate winner-takes-all
selection can cause oscillation. A stable selector therefore resembles a
decision process with hysteresis. It can require a candidate to win for several
consecutive observations and exceed the current motion by a margin before a
relevance-driven switch is accepted.

Stability alone can create repetitive behaviour. If the same catalogue item
remains marginally best, the robot may repeat it indefinitely even when several
alternatives are musically acceptable. Bounded-hold diversity provides a
separate mechanism: after a configured duration, selection may rotate among a
near-optimal set without pretending that the music itself has changed. This
separation between relevance and diversity is conceptually important. It
prevents diversity pressure from overriding a strong musical match while still
making a finite library appear less mechanical.

\section{Motion Transitions and Beat--Phase Synchronisation}
\label{sec:lit_transitions_phase}

\subsection{State-compatible motion entry}
\label{subsec:lit_state_entry}

Motion-graph and motion-matching research establishes pose similarity as a
foundation for transition search \cite{kovar2023motion,holden2020learned}.
For humanoid dance, pose distance is necessary but incomplete. Two frames can
have similar joint angles while their velocities point in opposite directions.
Likewise, geometrically similar poses can place different feet in support,
causing visible sliding or an abrupt change in the root trajectory.

A more complete entry cost for candidate phase $\phi$ is

\begin{equation}
\begin{split}
    C(\phi) ={}&
      \lambda_q C_q(\phi)
      + \lambda_v C_v(\phi)
      + \lambda_c C_c(\phi) \\
      &+ \lambda_r C_r(\phi)
      + \lambda_s C_s(\phi),
\end{split}
\end{equation}

where the terms measure joint pose, joint velocity, contact state, root state,
and musical or key-pose salience. The weights expose a practical compromise. A
large pose weight favours a visually smooth first frame; a large salience weight
favours musically meaningful entry; contact and root terms reduce foot and
world-frame artefacts. The transition planner can restrict its search to
salient phases, thereby reducing computation and avoiding entry in the middle
of an expressive gesture.

\subsection{Blending, root alignment, and output constraints}
\label{subsec:lit_blending}

Once an entry state is chosen, blending can distribute the remaining mismatch
over time. Database animation methods use automatically generated transitions
to connect captured sequences \cite{kovar2023motion}. For a floating-base
humanoid, the target motion should first be aligned to the current world root;
otherwise a locally smooth joint blend can still cause a global translation or
yaw jump. Joint interpolation should also respect angular topology and use the
shortest path for wrapped angles.

The blend duration is not purely aesthetic. If the pose displacement for joint
$j$ is $\Delta q_j$, then a lower bound derived from a velocity limit
$\dot{q}^{\max}_j$ is approximately

\begin{equation}
    T_{\mathrm{blend}}
    \geq
    \max_j \frac{|\Delta q_j|}{\dot{q}^{\max}_j}.
\end{equation}

Acceleration limits can impose a longer bound. Quantising the final duration to
a beat subdivision preserves a clearer musical structure, but can increase
scheduling latency. A final limiter remains necessary because phase-speed
adaptation, interpolation, and numerical sampling can combine to produce a
larger step than any individual component predicts.

These mechanisms improve kinematic continuity but do not prove dynamic
stability. A smooth joint-space path may still require infeasible torques or
move the centre of mass outside a viable support region. This distinction is
central when interpreting MuJoCo output later in the dissertation.

\subsection{Beat-to-key-pose alignment}
\label{subsec:lit_beat_keypose}

Music-conditioned generators often measure whether kinetic events coincide
with musical beats, and Beat-It makes beats and key poses explicit conditioning
variables \cite{huang2024beat}. A retrieval system faces the inverse problem:
the motion already contains salient events, so playback timing must adapt those
events to the incoming beat stream.

Instantaneously setting motion phase at every beat guarantees small phase error
at the update instant but creates a discontinuity in motion time. A gradual
controller instead modifies phase rate and consumes error over an interval. If
$\phi_t$ is the current cyclic phase, $\hat{\phi}_t$ is the desired key-pose
phase, and $\omega_t$ is playback rate, a simplified update is

\begin{equation}
    \phi_{t+1}
    = \phi_t + \Delta t\,
      \operatorname{clip}\!\left(
        \omega_t + k_p\,\operatorname{wrap}(\hat{\phi}_t-\phi_t),
        \omega_{\min},\omega_{\max}
      \right).
\end{equation}

The bounded rate prevents an accurate beat estimate from demanding an
unacceptable motion speed. The controller must also distinguish beat tracking
error from structural incompatibility. If a motion has too many or too few
salient poses for the observed rhythm, phase correction alone cannot solve the
problem; the motion should receive a lower compatibility score during
selection.

\section{Humanoid Motion Retargeting and Validation}
\label{sec:lit_retargeting}

\subsection{The embodiment gap}
\label{subsec:lit_embodiment_gap}

Retargeting maps a source performance to a target body while attempting to
preserve its semantic and perceptual character. For humanoid robots, this task
is constrained by differences in morphology and mechanical capability. SMPL
provides a consistent human source representation \cite{loper2023smpl}, but a
Unitree G1 trajectory requires a target-specific floating base and 29 actuated
joint coordinates. Direct copying is impossible because joint meanings, axes,
and ranges do not correspond one to one.

GMR treats motion retargeting as a consequential part of the humanoid learning
pipeline rather than a neutral preprocessing operation
\cite{araujo2025retargeting}. Its evaluation shows that artefacts such as foot
sliding, self-penetration, and physically implausible reference motion can
reduce downstream tracking robustness. This observation supports two design
decisions in the Humanoid Matcher. First, retargeting occurs offline so that its
cost and failures do not enter the live audio path. Second, the catalogue stores
retargeting provenance and admits only trajectories that satisfy explicit
preflight checks.

Retargeting quality has at least three dimensions. \emph{Fidelity} measures how
well the target preserves the source gesture. \emph{Kinematic validity} covers
joint ranges, continuity, grounding, and collision geometry. \emph{Dynamic
trackability} concerns whether a feedback controller and physical robot can
execute the reference under contact and torque constraints. A trajectory can
pass the first two dimensions and still fail the third. The experiments in this
dissertation should therefore avoid using kinematic success as evidence of
physical deployment readiness.

\subsection{Simulation and the limits of kinematic preview}
\label{subsec:lit_simulation}

MuJoCo is designed for multi-joint dynamics and contact-rich model-based control
\cite{todorov2012mujoco}. The capabilities of the simulator, however, do not
determine how a particular application uses it. Writing a target pose directly
to the simulator state and invoking forward kinematics visualises the robot
configuration and allows geometric checks, but it bypasses actuator dynamics,
balance control, and disturbance response.

This distinction prevents an otherwise common category error. A motion displayed
in a physics engine is not necessarily a physics-controlled motion. Kinematic
preview remains valuable: it can reveal joint-order mistakes, orientation
errors, ground penetration, discontinuities, and self-collision before an
online candidate is accepted. It should be reported as preflight and
visualisation evidence, not as a sim-to-real result.

\section{Critical Synthesis and Research Gap}
\label{sec:lit_synthesis}

The reviewed literature provides strong solutions to individual components of
music-driven motion, but no single line of work directly satisfies all of the
constraints of the present project.

\begin{description}
    \item[Offline dance generation]
    FACT, Bailando, EDGE, FineDance, and Beat-It provide increasingly expressive
    mappings from music to human motion
    \cite{li2021ai,siyao2022bailando,tseng2023edge,li2023finedance,huang2024beat}.
    Their strengths are learned choreography, diversity, and long-range
    structure. Their typical output is not automatically a target-robot action,
    and offline results do not establish strict causal latency.

    \item[Streaming generation]
    DiscoForcing addresses causality, bounded latency, and generated-history
    stability directly \cite{ji2026discoforcing}. It offers a compelling route
    to novel online motion, but requires a trained generative model and a
    deployment stack able to manage its output distribution.

    \item[Direct humanoid audio control]
    RoboPerform reduces the human-to-robot pipeline by learning direct
    audio-conditioned humanoid actions \cite{li2026you}. This strengthens
    embodiment integration, while moving the burden toward large-scale policy
    learning, training data, and sim-to-real validation.

    \item[Database animation]
    Motion graphs and motion matching preserve captured detail and support
    responsive reuse of motion data
    \cite{kovar2023motion,holden2020learned}. They provide principles for
    state-compatible transition search but do not by themselves define a music
    query or beat-aware switching policy.

    \item[Humanoid retargeting]
    GMR addresses the embodiment gap and demonstrates that reference quality
    affects downstream tracking \cite{araujo2025retargeting}. Retargeting alone
    does not decide which trajectory fits live music or how successive
    trajectories should be scheduled.
\end{description}

The resulting research gap is an integrated, lightweight system that treats
music similarity, robot-motion compatibility, temporal selection stability,
beat alignment, transition continuity, and catalogue validity as separate but
cooperating problems. A straightforward nearest-neighbour system is
insufficient because it can select tempo-incompatible motion and oscillate
between near-equal candidates. A beat-only controller is insufficient because
tempo does not encode style or activity. A smooth blend is insufficient if the
entry states disagree in velocity or support. Finally, a retargeted motion is
insufficient if its provenance and preflight status are unknown.

The Humanoid Matcher is motivated by this gap. It combines a rolling audio
representation with hierarchical retrieval, explicitly interpretable
compatibility terms, stateful switch acceptance, beat-boundary scheduling,
key-pose phase control, state-aware entry search, root-aligned blending, and
offline G1 validation. Its intended contribution is therefore architectural
and experimental rather than a new foundational generative model. The relevant
question is whether these components together achieve a useful operating point:
less novel than unconstrained generation, but causal, inspectable,
computationally modest, and bounded by a prevalidated motion library.

\section{Chapter Summary}
\label{sec:lit_summary}

This chapter reviewed the foundations of real-time music-driven humanoid
motion. Music information retrieval provides complementary temporal and
semantic representations, while causal beat analysis exposes the need to
balance local evidence against rhythmic consistency. AIST++ and FineDance
support large-scale study of music--dance relationships, but their human motion
representations must be distinguished from executable robot coordinates.
Autoregressive, discrete, and diffusion models demonstrate increasingly
expressive choreography; RoboPerform and DiscoForcing extend the field toward
direct humanoid action and strict streaming generation.

Example-based animation contributes a different set of principles: preserve
high-quality captured motion, search explicitly over a database, and join clips
only at compatible states. GMR further shows that retargeting quality materially
affects humanoid tracking, supporting offline retargeting and preflight before a
motion enters the online catalogue. Together, these findings justify the hybrid
retrieval-and-control approach developed in the following chapters. Chapter~3
translates the identified requirements into the offline and online architecture
of the Humanoid Matcher.
